\documentclass[11pt,a4paper]{report}
\usepackage[portuguese]{babel}
\usepackage[margin=1in]{geometry}
\usepackage{graphicx}
\usepackage{xcolor}
\usepackage{listings}
\usepackage{fancyhdr}
\usepackage{hyperref}
\usepackage{amsmath}
\usepackage{float}
\usepackage{array}
\usepackage{tabularx}

% Code highlighting
\lstset{
    language=Python,
    basicstyle=\ttfamily\small,
    keywordstyle=\color{blue}\bfseries,
    commentstyle=\color{gray},
    stringstyle=\color{red},
    breaklines=true,
    showstringspaces=false,
    tabsize=4,
    frame=single,
    numbers=left,
    numberstyle=\tiny\color{gray}
}

% Header and footer
\pagestyle{fancy}
\fancyhf{}
\rhead{Processamento de Linguagens}
\lhead{TP - Linguagem LFun}
\cfoot{\thepage}

\title{\textbf{Relatório do Trabalho Prático} \\[0.5cm]
        \large Processamento de Linguagens \\[0.3cm]
        \large Implementação de Analisador e Interpretador para a Linguagem LFun}
\author{LESI - 2º Ano, 2º Semestre}
\date{\today}

\begin{document}

\maketitle

\newpage
\tableofcontents
\newpage

% ==================== INTRODUÇÃO ====================
\chapter{Introdução}

Este relatório documenta o trabalho prático da unidade curricular de \textbf{Processamento de Linguagens}, desenvolvido no contexto da Licenciatura em Engenharia de Sistemas Informáticos (LESI), 2º ano, 2º semestre.

O trabalho centrou-se no desenvolvimento progressivo de um analisador completo para uma linguagem funcional denominada \textbf{LFun} (Language Functional), incluindo componentes de análise léxica, análise sintática, avaliação básica e, finalmente, um sistema robusto com representação intermédia e suporte avançado a padrões.

O projeto foi estruturado em quatro fases incrementais (A, B, C e D), permitindo uma construção gradual e consolidação de conceitos fundamentais no processamento de linguagens.

% ==================== OBJECTIVOS ====================
\chapter{Objectivos}

O presente trabalho teve como objectivos principais:

\begin{enumerate}
    \item Compreender e aplicar os conceitos teóricos de processamento de linguagens, nomeadamente:
    \begin{itemize}
        \item Análise léxica e tokenização
        \item Análise sintática e gramáticas formais
        \item Representação intermédia (IR)
        \item Avaliação e interpretação de código
    \end{itemize}
    
    \item Implementar um analisador léxico capaz de:
    \begin{itemize}
        \item Identificar palavras-chave, operadores e literais
        \item Processar comentários de linha e bloco
        \item Gerar tokens correctamente
    \end{itemize}
    
    \item Desenvolver um analisador sintático que:
    \begin{itemize}
        \item Verifique a conformidade estrutural do código
        \item Valide expressões complexas
        \item Produza uma representação sintáctica válida
    \end{itemize}
    
    \item Criar um interpretador com funcionalidades avançadas:
    \begin{itemize}
        \item Avaliação de expressões
        \item Suporte a pattern matching robusto
        \item Representação intermédia genérica (AST)
    \end{itemize}
    
    \item Garantir qualidade através de testes abrangentes
\end{enumerate}

% ==================== ESTRUTURA DO PROJECTO ====================
\chapter{Estrutura do Projeto}

\section{Organização Geral}

O projeto foi desenvolvido em quatro etapas, cada uma construindo sobre a anterior:

\begin{center}
\begin{tabularx}{\textwidth}{|p{2cm}|X|p{6cm}|}
\hline
\textbf{Parte} & \textbf{Objectivo} & \textbf{Componentes} \\
\hline
\textbf{A} & Análise Léxica & lexer.py, lexer\_test \\
\hline
\textbf{B} & Análise Sintática & lexer.py, grammar.py, parser \\
\hline
\textbf{C} & Avaliação Básica & parser.py, main.py, test \\
\hline
\textbf{D} & IR e Pattern Matching & ast\_nodes.py, eval.py, test\_eval.py \\
\hline
\end{tabularx}
\end{center}

\section{Arquitectura do Sistema}

O sistema segue uma arquitetura em camadas:
\begin{figure}[H]
\centering
\begin{lstlisting}
┌─────────────────────────────────────────┐
│      Código Fonte (Entrada)             │
└─────────────────┬───────────────────────┘
        ┌─────────▼────────────┐
        │  ANÁLISE LÉXICA      │
        │  (Lexer)             │
        │  lexer.py            │
        └─────────┬────────────┘
        ┌─────────▼────────────┐
        │  ANÁLISE SINTÁTICA   │
        │  (Parser)            │
        │  parser.py           │
        └─────────┬────────────┘
        ┌─────────▼──────────────────────┐
        │  REPRESENTAÇÃO INTERMÉDIA      │
        │  (AST - Árvore de Sintaxe      │
        │   Abstracta)                   │
        │  ast_nodes.py                  │
        └─────────┬──────────────────────┘
        ┌─────────▼────────────┐
        │  AVALIAÇÃO           │
        │  (Evaluator)         │
        │  eval.py             │
        └─────────┬────────────┘
        ┌─────────▼───────────────────────┐
        │  Resultado (Saída)              │
        └─────────────────────────────────┘
\end{lstlisting}
\end{figure}
\subsection{Componentes Principais}

\subsubsection{Análise Léxica}
O módulo \texttt{lexer.py} implementa a tokenização do código fonte, identificando:
\begin{itemize}
    \item Palavras-chave (let, fun, if, match, etc.)
    \item Operadores (+, -, *, /, ==, <>, etc.)
    \item Literais numéricos e booleanos
    \item Identificadores
    \item Comentários


\end{itemize}

\subsubsection{Análise Sintática}
O módulo \texttt{parser.py} utiliza a biblioteca PLY (Python Lex-Yacc) para:
\begin{itemize}
    \item Processar tokens
    \item Validar a gramática
    \item Construir uma Árvore de Sintaxe Abstracta (AST)
\end{itemize}

\subsubsection{Representação Intermédia}
O ficheiro \texttt{ast\_nodes.py} define classes para todas as construções:
\begin{itemize}
    \item \texttt{NumberLiteral}, \texttt{BoolLiteral}
    \item \texttt{VarExpr}, \texttt{IfExpr}, \texttt{MatchExpr}
    \item \texttt{BinOp}, \texttt{UnaryOp}, \texttt{CallExpr}
    \item Padrões: \texttt{IntPattern}, \texttt{WildcardPattern}, \texttt{VarPattern}
    \item Definições: \texttt{LetStmt}, \texttt{FunDef}
\end{itemize}

\subsubsection{Avaliação}
O módulo \texttt{eval.py} implementa a interpretação de programas:
\begin{itemize}
    \item Avaliação recursiva de expressões
    \item Pattern matching robusto
    \item Gestão de contexto e variáveis
    \item Suporte a funções
\end{itemize}

% ==================== IMPLEMENTAÇÃO ====================
\chapter{Implementação}

\section{Parte A - Análise Léxica}

A primeira etapa focou-se na construção de um analisador léxico robusto.

\subsection{Funcionalidades}

\begin{itemize}
    \item Extracção de tokens do código fonte
    \item Identificação de palavras reservadas
    \item Reconhecimento de operadores e delimitadores
    \item Processamento de comentários de linha (\texttt{--}) e bloco (\texttt{\{- -\}})
    \item Suporte a literais numéricos e booleanos
\end{itemize}

\subsection{Estrutura do Lexer}

\begin{lstlisting}
# Exemplo de tokenização
Entrada: let x = 42 
Tokens:  ['let', 'x', '=', '42']
\end{lstlisting}

\vspace{250}
\section{Parte B - Análise Sintática}

A segunda etapa focou-se na construção da gramática e parser.

\subsection{Gramática}

A gramática define a sintaxe válida para a linguagem LFun:

\begin{lstlisting}
program : expr
        | stmt_list

expr : NUMBER
     | TRUE | FALSE
     | ID
     | expr '+' expr
     | expr '-' expr
     | expr '*' expr
     | expr '/' expr
     | 'if' expr 'then' expr 'else' expr
     | 'match' expr 'with' case_list
     | ID '(' expr ')'
\end{lstlisting}

\subsection{Utilização de PLY}

O projecto utiliza a biblioteca \textbf{PLY} (Python Lex-Yacc) para:
\begin{itemize}
    \item Definição declarativa da gramática
    \item Geração automática do parser
    \item Tratamento de conflitos de ambiguidade
    \item Recuperação de erros de sintaxe
\end{itemize}
\vspace{270}
\section{Parte C - Avaliação Básica}

A terceira etapa integrou o parser com um sistema básico de avaliação.

\subsection{Características}

\begin{itemize}
    \item Interpretação de expressões aritméticas
    \item Avaliação de expressões condicionais (if-then-else)
    \item Processamento de padrões simples
    \item Interface REPL (Read-Eval-Print Loop)
\end{itemize}

\section{Parte D - Representação Intermédia e Pattern Matching}

A quarta e final etapa implementou um sistema robusto com AST e pattern matching avançado.

\subsection{Nós da AST}

\begin{lstlisting}
class NumberLiteral(Node):
    def __init__(self, value):
        self.value = int(value)

class BoolLiteral(Node):
    def __init__(self, value):
        self.value = bool(value)

class VarExpr(Node):
    def __init__(self, name):
        self.name = name

class BinOp(Node):
    def __init__(self, op, left, right):
        self.op = op
        self.left = left
        self.right = right

class MatchExpr(Node):
    def __init__(self, expr, cases):
        self.expr = expr
        self.cases = cases

class Case(Node):
    def __init__(self, patterns, expr):
        self.patterns = patterns
        self.expr = expr
\end{lstlisting}
\vspace{70}
\subsection{Pattern Matching}

O sistema suporta múltiplos tipos de padrões:

\begin{itemize}
    \item \textbf{Padrões Literais:} Correspondência exacta de valores
    \item \textbf{Alternativas:} Múltiplos padrões com o operador |
    \item \textbf{Wildcard:} Símbolo \texttt{\_} para qualquer valor
    \item \textbf{Captura de Variáveis:} Nomes começados com minúscula
\end{itemize}

\subsection{Exemplos de Pattern Matching}

\begin{lstlisting}
# Padrão com alternativas
match x with
| 1 | 2 | 3 -> "um, dois ou tres"
| _ -> "outro valor"

# Captura de variável
match x with
| n when n > 0 -> "positivo"
| n -> n + 1

# Padrão literal com wildcard
match (x, y) with
| (0, _) -> "x é zero"
| (_, 0) -> "y é zero"
| _ -> "ambos não-zero"
\end{lstlisting}

\subsection{Avaliação Genérica}

O avaliador opera exclusivamente sobre nós AST:

\begin{lstlisting}
def evaluate(node, context):
    if isinstance(node, NumberLiteral):
        return node.value
    elif isinstance(node, BoolLiteral):
        return node.value
    elif isinstance(node, VarExpr):
        return context.get(node.name)
    elif isinstance(node, BinOp):
        left = evaluate(node.left, context)
        right = evaluate(node.right, context)
        return apply_op(node.op, left, right)
    elif isinstance(node, MatchExpr):
        return evaluate_match(node, context)
    # ... outros casos
\end{lstlisting}

% ==================== TECNOLOGIAS ====================
\chapter{Tecnologias Utilizadas}

\begin{itemize}
    \item \textbf{Python 3:} Linguagem de implementação
    \item \textbf{PLY (Python Lex-Yacc):} Framework para lexer e parser
    \item \textbf{Pytest:} Framework para testes
\end{itemize}

% ==================== TESTES ====================
\chapter{Testes e Validação}

\section{Estratégia de Testes}

O projecto inclui múltiplos níveis de testes:

\subsection{Testes Unitários}

Testes dos componentes individuais através de \texttt{test\_eval.py}:

\begin{lstlisting}
# Testes básicos
def test_absolute():
    assert evaluate("absoluto(-5)") == 5
    assert evaluate("absoluto(5)") == 5

def test_boolean():
    assert evaluate("eZero(0)") == True
    assert evaluate("eZero(1)") == False

def test_alternatives():
    assert evaluate("avalia(-1)") == True
    assert evaluate("avalia(0)") == False
    assert evaluate("avalia(2)") == True
\end{lstlisting}

\subsection{Cobertura de Testes}

Os testes cobrem:
\begin{itemize}
    \item Expressões aritméticas
    \item Expressões booleanas
    \item Padrões literais
    \item Padrões com alternativas
    \item Wildcard patterns
    \item Captura de variáveis
    \item Nested match
    \item Casos não-exaustivos
\end{itemize}

\section{Resultados de Teste}

Exemplos de execução bem-sucedida:

\begin{lstlisting}
absoluto(-5) = 5
absoluto(5) = 5
eZero(0) = True
eZero(1) = False
avalia(-1) = True
avalia(0) = False
avalia(2) = True
sign(-5) = -1
sign(0) = 0
sign(7) = 1
\end{lstlisting}

% ==================== RESULTADOS ====================
\chapter{Resultados e Demonstração}

\section{Interface REPL}

O projeto inclui um REPL (Read-Eval-Print Loop) interactivo:

\begin{lstlisting}
$ python D/main.py
LFun Interpreter 1.0
Type 'exit' or 'quit' to exit

> let absoluto x = match x with | n when n < 0 -> -n | n -> n
> absoluto(-10)
10

> match 5 with | 1 | 2 | 3 -> "pequeno" | _ -> "grande"
"grande"
\end{lstlisting}

\section{Exemplo Completo}

Um exemplo demonstrando vários aspetos da linguagem:

\begin{lstlisting}
-- Definição de função com pattern matching
let sign x = match x with
| 0 -> 0
| n when n < 0 -> -1
| _ -> 1

-- Utilização da função
sign(-5)    -- retorna: -1
sign(0)     -- retorna: 0
sign(7)     -- retorna: 1

-- Pattern matching com alternativas
match 2 with
| 1 | 2 | 3 -> "um dois ou três"
| _ -> "outro"
-- retorna: "um dois ou três"

-- Expressões condicionais
if 5 > 3 then "verdadeiro" else "falso"
-- retorna: "verdadeiro"
\end{lstlisting}

\section{Conformidade com Requisitos}

\begin{table}[H]
\centering
\begin{tabularx}{\textwidth}{|p{3cm}|p{2cm}|X|}
\hline
\textbf{Requisito} & \textbf{Status} & \textbf{Evidência} \\
\hline
Pattern Matching (literal, alternativas, wildcard, captura) & ✓ & test\_eval\_more.py \\
\hline
Representação Intermédia Genérica & ✓ & ast\_nodes.py \\
\hline
Avaliação sobre IR & ✓ & eval.py \\
\hline
Suporte a Funções & ✓ & let, fun, signatures \\
\hline
Testes Abrangentes & ✓ & test\_eval.py, test/ \\
\hline
Qualidade de Código & ✓ & Sem código duplicado \\
\hline
\end{tabularx}
\end{table}

% ==================== CONCLUSÃO ====================
\chapter{Conclusão}

Este trabalho prático permitiu consolidar conhecimentos fundamentais sobre processamento de linguagens, desde a análise léxica até à interpretação avançada com representação intermédia.

\section{Aprendizagens Principais}

\begin{enumerate}
    \item \textbf{Separação de Responsabilidades:} Cada etapa (lexer, parser, AST, avaliador) tem um papel bem definido
    \item \textbf{Representação Intermédia:} Decisão crítica de usar AST em vez de dicionários permite maior flexibilidade
    \item \textbf{Pattern Matching:} Implementação robusto de múltiplos tipos de padrões
    \item \textbf{Testes:} Importância de cobertura abrangente para validar comportamento
\end{enumerate}

\section{Funcionalidades Implementadas}

\begin{itemize}
    \item Análise léxica completa com suporte a comentários
    \item Parser robusto com tratamento de ambiguidades
    \item Representação Intermédia genérica (AST)
    \item Avaliador recursivo com pattern matching avançado
    \item Suporte a funções e assinaturas de tipo
    \item REPL interactivo
    \item Suite de testes abrangente
\end{itemize}

\section{Possíveis Melhorias Futuras}

\begin{itemize}
    \item Optimizações de performance (compilação ou caching)
    \item Geração de código intermediário (bytecode)
    \item Tratamento de erros mais informativo
\end{itemize}

\section{Reflexão Final}

O projecto foi realizado com sucesso, cumprindo todos os requisitos especificados. O sistema implementado é robusto, bem-estruturado e facilmente extensível. A arquitectura em camadas permite fáceis mudanças futuras sem afectar outras partes do sistema.

A escolha de utilizar uma AST genérica desde o início mostrou-se acertada, permitindo operações de transformação e análise sofisticadas no futuro.

\newpage
\appendix

\chapter{Estrutura de Ficheiros}

\begin{lstlisting}
TP_PL202526/
├── README.md
├── pl-tp-enunciado_pub.pdf
├── A/
│   ├── lexer.py
│   └── lexer_test
├── B/
│   ├── lexer.py
│   ├── grammar.py
│   ├── parser.py
│   └── parsetab.py
├── C/
│   ├── lexer.py
│   ├── parser.py
│   ├── main.py
│   ├── parsetab.py
│   └── test
├── D/
│   ├── ast_nodes.py
│   ├── eval.py
│   ├── grammar.py
│   ├── lexer.py
│   ├── main.py
│   ├── README.md
│   ├── test_eval.py
│   ├── test_eval_more.py
│   ├── run_tests.py
│   └── test/
│       └── test_eval_pytest.py
\end{lstlisting}

\end{document}
